\section*{Chapter \ref{ch:m2:fixings-and-flows} --- Fixings and Flows}

\begin{solution}{exo:m2:fixings-and-flows:1}
Its index is computed with currencies valued at the fix. A trade at the fix
matches the index exactly; a trade earlier, even at a better price, leaves a
difference against the index that the fund must explain as tracking error.
\end{solution}

\begin{solution}{exo:m2:fixings-and-flows:2}
Its hedge must grow by $0.7 \times 3\% \times 50$ billion: it sells USD~1.05
billion forward against yen.
\end{solution}

\begin{solution}{exo:m2:fixings-and-flows:3}
On 15 February 2015 WM widened the window from one to five minutes, from 2.5
minutes before to 2.5 minutes after the hour, and added price sources. To move
the average price of five minutes a trader must trade much more, for longer and
more visibly than to move one minute.
\end{solution}

\begin{solution}{exo:m2:fixings-and-flows:4}
The fix is $3 \times 1.3/2 = 1.95$ pips above the start; the dealer's cost is still
1.5, so it earns 0.45 pips on EUR~1 billion, USD~45\,000.
\end{solution}

\begin{solution}{exo:m2:fixings-and-flows:5}
It buys USD~1.2 billion of American assets and sells USD~1.2 billion of European
ones: the losing region is bought back to its weight, and dollars are bought.
\end{solution}

\begin{solution}{exo:m2:fixings-and-flows:6}
Long the option, it is long gamma: it sells EURUSD as the price rises towards the
strike and buys as it falls, damping moves and pinning the price near the strike
into the cut. At the cut the delta hedge must be unwound or completed.
\end{solution}

\begin{solution}{exo:m2:fixings-and-flows:7}
At $p = 0.5$, a mean of about USD~74\,346 and a standard deviation of about
USD~100\,276; at $p = 1$, about USD~148\,692 and USD~200\,552.
\end{solution}

\begin{solution}{exo:m2:fixings-and-flows:8}
The dealer's cost is the same, but the fix is not: buying before the window lets
the dealer's own impact raise the fix at which it sells to the client. The client
pays $\lambda Q p/2$ per unit more than with no pre-hedging, and that is the
dealer's profit. Whether pre-hedging is legitimate depends on whether it is done
to manage the dealer's risk in a way that benefits the client, not on the cost.
\end{solution}

\begin{solution}{pb:m2:fixings-and-flows:1}
\textbf{1.} $0.3 \times 10 = 3$ pips.
\textbf{2.} 1.5 pips above the start.
\textbf{3.} Also 1.5 pips; profit zero.
\textbf{4.} None from impact in this model; in reality the uncertainty of prices
within the window against the fix, and the operational risk of executing a
billion in five minutes.
\textbf{5.} A fee or spread agreed in advance, as a few pips added to the fix,
or a commission.
\textbf{6.} 2.25 pips above the start.
\textbf{7.} USD~75\,000, paid by the client through a fix 0.75 pips higher.
\textbf{8.} About USD~100\,000, from the random move between the pre-hedge and
the window.
\textbf{9.} USD~150\,000 expected, with a standard deviation of about USD~200\,000.
\textbf{10.} Every unit bought raises the price by the same amount, whenever it is
bought; the order of purchases changes who pays what, not the total cost.
\textbf{11.} The dealer's net need falls to EUR~400 million if it can match the
two orders internally at the fix, and the impact with it; matching at the fix is
fair to both clients.
\textbf{12.} Dollar sales push EURUSD up; around month end, and some of it before
the window if others trade ahead.
\textbf{13.} Because those who expect the flow trade before it, and the dealer's
own pre-hedging moves the price before the window.
\textbf{14.} By splitting the order over a longer period or an algorithm, by
asking dealers for a price earlier, or by negotiating how the dealer executes and
what it charges.
\textbf{15.} Not to share the order or collude, not to influence the fix to
benefit from it, to price transparently, and to pre-hedge only as principal, in a
way designed to benefit the client, and after telling it how it does so.
\textbf{16.} When the dealer bears a real risk from the anticipated order, trades
to reduce it in a way meant to benefit the client, and has told the client.
\textbf{17.} Collusion: traders at competing banks shared clients' orders and
coordinated their trading, and withheld bids or offers, to move the fixes; the
US charge was a conspiracy to fix prices and rig bids.
\textbf{18.} Moving a five-minute average takes more trading, for longer, and is
easier to see; it does not stop a large order from moving the price, nor a dealer
from trading ahead of it.
\textbf{19.} Named result: \emph{the \omterm{def:m2:fixings-and-flows:fix}{fixing order}'s profit}, zero without
pre-hedging, USD~75\,000 with half pre-hedged and USD~150\,000 with all
pre-hedged, on EUR~1 billion, paid by the client through a higher fix.
\textbf{20.} The client, through a fix moved by the dealer's own trades.
\end{solution}

\begin{solution}{iq:m2:fixings-and-flows:1}
A benchmark exchange rate published for each pair at 4pm London time from trades
and quotes in a five-minute window (one minute before February 2015). Asset
managers and index providers value portfolios and indices with it, and many
clients' orders and derivatives are priced at it.
\iqlookfor{the window, and why managers trade at it.}
\end{solution}

\begin{solution}{iq:m2:fixings-and-flows:2}
They follow the month's asset returns: hedged foreign investors adjust hedges to
the new value of their holdings, and fixed-weight portfolios sell what rose and
buy what fell. After a strong month for US equities, foreign hedgers sell more
dollars, and global funds sell US assets and dollars: dollar sales.
\iqlookfor{the two mechanisms and the sign.}
\end{solution}

\begin{solution}{iq:m2:fixings-and-flows:3}
Net it against other \omterm{def:m2:fixings-and-flows:fix}{fixing orders} and flows; execute the remainder mostly in the
window, evenly or with an algorithm matching the window's weighting; pre-hedge
only if the risk warrants it, as principal and in a way meant to benefit the
client; tell the client beforehand how the order will be handled and what it
will cost, and document the execution.
\iqlookfor{netting, window execution, disclosure and the Code.}
\end{solution}

\begin{solution}{iq:m2:fixings-and-flows:4}
With permanent impact $\lambda$, buying $pQ$ before the window at average
$P_0 + \lambda pQ/2$ and the rest in the window, whose average is the fix
$P_0 + \lambda Q(1+p)/2$, gives an average cost $P_0 + \lambda Q/2$ whatever $p$.
The profit $\lambda Q^2 p/2$ comes from the client, who pays a fix raised by the
dealer's own earlier buying.
\iqlookfor{the cost invariance and the source of the profit.}
\end{solution}

\begin{solution}{iq:m2:fixings-and-flows:5}
Traders at several banks shared clients' orders in chat rooms and coordinated
trading to move the 4pm and ECB fixes; regulators in the UK, US and elsewhere
fined banks billions in 2014, and four pleaded guilty to US criminal charges in
2015. The fix window was widened to five minutes, the \omterm{def:m2:fx-spot-microstructure:code}{FX Global Code} set
standards for \omterm{def:m2:fixings-and-flows:fix}{fixing orders} and pre-hedging, and banks tightened their controls.
\iqlookfor{facts, outcomes and reforms.}
\end{solution}

\begin{solution}{iq:m2:fixings-and-flows:6}
Record every order and trade with timestamps; for each trader and account,
measure the share of volume traded in and just before fix windows, the direction
relative to their \omterm{def:m2:fixings-and-flows:fix}{fixing orders} and positions, and the price impact in the
window against a model of normal activity; flag concentrations that move the fix
in the direction of a position, and link them with communications surveillance.
Test with synthetic cases and review alerts' precision.
\iqlookfor{features tied to positions and fixes, baselines, and communications.}
\end{solution}
